\begin{thebibliography}{99}
\footnotesize
\bibitem{kawata} Y. Kawata, «The Great Hanshin-Awaji Earthquake Disaster», \textit{J. Natural Disaster Science}, vol.~17, n.º~2, 1995. Las proporciones por día provienen de un trabajo del mismo autor de 1996, consultado a través de una cita secundaria.
\bibitem{noji} E. K. Noji, «Public health consequences of earthquakes», Supercourse, Universidad de Pittsburgh. [En línea]. Disponible: \url{https://sites.pitt.edu/~super1/lecture/lec13051/013.htm}
\bibitem{galvis2020} F. Galvis, E. Miranda, P. Heresi, H. Dávalos y J. Ruiz-García, «Overview of collapsed buildings in Mexico City after the 19 September 2017 (Mw7.1) earthquake», \textit{Earthquake Spectra}, 2020, doi: 10.1177/8755293020936694.
\bibitem{murphy2004} R. R. Murphy, «Trial by fire: activities of the rescue robots at the World Trade Center», \textit{IEEE Robot. Autom. Mag.}, vol.~11, n.º~3, pp.~50--61, 2004.
\bibitem{murphy2021} R. R. Murphy, «An expert on search and rescue robots explains the technologies used in disasters like the Florida condo collapse», \textit{The Conversation}, 2021.
\bibitem{tranngoc2023} P. T. Tran-Ngoc \textit{et al.}, «Intelligent insect--computer hybrid robot: installing innate obstacle negotiation and onboard human detection onto cyborg insect», \textit{Adv. Intell. Syst.}, vol.~5, n.º~5, 2200319, 2023.
\bibitem{lin2025} Q. Lin \textit{et al.}, «Cyborg insect factory: automatic assembly for insect-computer hybrid robot via vision-guided robotic arm manipulation of custom bipolar electrodes», \textit{Nat. Commun.}, vol.~16, 6073, 2025.
\bibitem{refat2025} C. M. M. Refat, M. Ariyanto, R. Tanaka, K. Yamamoto y K. Morishima, «Autonomous navigation of bio-intelligent cyborg insect based on insect visual perception», \textit{Adv. Intell. Syst.}, vol.~7, n.º~10, 2400838, 2025.
\bibitem{htx2025} Public Service Division, Gobierno de Singapur, «From lab to rubble», 12 ago. 2026; y «Singapore's cyborg cockroaches helping with search and rescue efforts in Myanmar quake», \textit{The Straits Times}, 5 abr. 2025.
\bibitem{latif2012} T. Latif y A. Bozkurt, «Line following terrestrial insect biobots», en \textit{Proc. IEEE EMBC}, 2012, pp.~972--975.
\bibitem{erickson2015} J. C. Erickson, M. Herrera, M. Bustamante, A. Shingiro y T. Bowen, «Effective stimulus parameters for directed locomotion in Madagascar hissing cockroach biobot», \textit{PLoS ONE}, vol.~10, n.º~8, e0134348, 2015.
\bibitem{sanchez2015} C. J. Sanchez, C.-W. Chiu, Y. Zhou, J. M. González, S. B. Vinson y H. Liang, «Locomotion control of hybrid cockroach robots», \textit{J. R. Soc. Interface}, vol.~12, 20141363, 2015.
\bibitem{ma2023} S. Ma, Y. Chen, S. Yang, S. Liu, L. Tang, B. Li y Y. Li, «The autonomous pipeline navigation of a cockroach bio-robot with enhanced walking stimuli», \textit{Cyborg Bionic Syst.}, vol.~4, 0067, 2023.
\bibitem{lin2023} Q. Lin \textit{et al.}, «Resilient conductive membrane synthesized by in-situ polymerisation for wearable non-invasive electronics on moving appendages of cyborg insect», \textit{npj Flex. Electron.}, vol.~7, 42, 2023.
\bibitem{bai2025} Y. Bai \textit{et al.}, «Swarm navigation of cyborg-insects in unknown obstructed soft terrain», \textit{Nat. Commun.}, vol.~16, 221, 2025.
\bibitem{gracidas2026} E. Gracidas-Reyes, G. Diaz-Arango, H. Vazquez-Leal y F. Marroquin-Gutierrez, «Temperature-dependent tethered locomotion behavior in the Madagascar hissing cockroach using a controlled-environment treadmill platform with exploratory illumination assays», \textit{Biology}, vol.~15, n.º~12, 947, 2026.
\bibitem{montagut2024} M. J. Montagut Marques \textit{et al.}, «Cyborg insect repeatable self-righting locomotion assistance using bio-inspired 3D printed artificial limb», \textit{npj Robot.}, vol.~2, 3, 2024.
\bibitem{dirafzoon2017} A. Dirafzoon, T. Latif, F. Gong, M. Sichitiu, A. Bozkurt y E. Lobaton, «Biobotic motion and behavior analysis in response to directional neurostimulation», en \textit{Proc. IEEE ICASSP}, 2017, pp.~2457--2461.
\bibitem{tranngoc2024} P. T. Tran-Ngoc, H. D. Nguyen, D. L. Le, R. Li, B. S. Chong y H. Sato, «Gait-adaptive navigation and human searching in field with cyborg insect», arXiv:2406.03713, 2024.
\bibitem{peterson1961} W. W. Peterson y D. T. Brown, «Cyclic codes for error detection», \textit{Proc. IRE}, vol.~49, n.º~1, pp.~228--235, 1961.
\bibitem{goertzel1958} G. Goertzel, «An algorithm for the evaluation of finite trigonometric series», \textit{Amer. Math. Monthly}, vol.~65, n.º~1, pp.~34--35, 1958.
\bibitem{volder1959} J. E. Volder, «The CORDIC trigonometric computing technique», \textit{IRE Trans. Electron. Comput.}, vol.~EC-8, n.º~3, pp.~330--334, 1959.
\bibitem{andraka1998} R. Andraka, «A survey of CORDIC algorithms for FPGA based computers», en \textit{Proc. ACM/SIGDA Int. Symp. Field Programmable Gate Arrays}, 1998, pp.~191--200.
\bibitem{librelane} FOSSi Foundation, «LibreLane». [En línea]. Disponible: \url{https://github.com/librelane/librelane} (consultado el 21 nov. 2026).
\bibitem{shalan2020} M. Shalan y T. Edwards, «Building OpenLANE: a 130nm OpenROAD-based tapeout-proven flow», en \textit{Proc. IEEE/ACM Int. Conf. Computer-Aided Design (ICCAD)}, 2020, pp.~1--6.
\bibitem{skywater} SkyWater Technology y Google, «SKY130 process design kit», biblioteca \texttt{sky130\_fd\_sc\_hd}. [En línea]. Disponible: \url{https://github.com/google/skywater-pdk}
\bibitem{lgvs} «Ley General de Vida Silvestre», última reforma publicada en el \textit{Diario Oficial de la Federación} el 16 jul. 2025.
\bibitem{gibbons2022} M. Gibbons, A. Crump, M. Barrett, S. Sarlak, J. Birch y L. Chittka, «Can insects feel pain? A review of the neural and behavioural evidence», \textit{Adv. Insect Physiol.}, pp.~155--229, 2022.
\bibitem{adamo2016} S. A. Adamo, «Do insects feel pain? A question at the intersection of animal behaviour, philosophy and robotics», \textit{Anim. Behav.}, vol.~118, pp.~75--79, 2016.
\bibitem{tinytapeout} Tiny Tapeout. [En línea]. Disponible: \url{https://tinytapeout.com} (consultado el 9 nov. 2026).
\end{thebibliography}
